% ST-NDT baseline for Section 4.2 (rebuttal change 1).
% Input from appendix/additional_results/additional_churchland_results.tex.
For ST-NDT (revision \texttt{92ffe21c2042eb886389b7ebc97d300c2fd23cac}), we use the authors' model and masking code unchanged, starting from their MC\_Maze configuration: four spatiotemporal Transformer layers with feedforward width 128, pre-normalization, a linear input embedding, and a linear readout of Poisson log-rates. The Transformer's width equals the number of input channels, so its representation has 191 dimensions instead of 192, and uses one attention head because 191 is prime. Inputs are raw channel counts, and attention spans all ten bins of a window. Training follows the published objective: Poisson prediction of masked counts (25\% of entries, set to zero) plus a weighted InfoNCE loss (temperature 0.07) between two views of each window in which 5\% of entries are zeroed or randomized. We use AdamW (weight decay $5\times10^{-5}$) with linear warm-up over the first 10\% of updates followed by cosine decay, for ten shuffled epochs (5,100 updates, $B=256$). The 12-run sweep crosses contrastive-loss weights $\{0.01,0.1,1\}$, peak learning rates $\{0.001,0.003\}$, and dropout rates $\{0.1,0.3\}$; it selects weight 1, learning rate 0.003, and dropout 0.3. Candidate latents are the coordinates of the final Transformer layer's output at the last bin of each window, the features that the authors' code exports for downstream analyses; because each window ends at its scored bin, they have no access to later activity. When fit to the training trials only, the selected setting predicts masked counts on held-out trials 0.123 bits per spike better than each channel's mean rate. ST-NDT sweep runs are ranked, and its coordinates are sign-oriented, thresholded, selected, and decoded, in the same way as for CEBRA and LangevinFlow.
